% =========================================================
% Chapter 6: Circular Motion and Applications of Newton's Laws
% POSN 1 Physics Preparation Notes
% Language: Thai
% Compiler: XeLaTeX
% =========================================================

\chapter{การเคลื่อนที่เป็นวงกลมและการประยุกต์กฎของนิวตัน}

\begin{center}
    {\Large \textbf{Circular Motion and Applications of Newton's Laws}}\\[4pt]
    {\normalsize เอกสารเตรียมสอบคัดเลือก สอวน. ฟิสิกส์ ค่าย 1}\\[6pt]
    {\normalsize จัดทำโดย \textbf{Tames Thanakrit Damduan}}\\[2pt]
    {\footnotesize \textcopyright\ 2026 Tames Thanakrit Damduan. All rights reserved.}\\[8pt]
\end{center}

\section*{เป้าหมายของบทเรียน}

หลังจากเรียนบทนี้ นักเรียนควรทำสิ่งต่อไปนี้ได้

\begin{enumerate}
    \item เข้าใจว่าอัตราเร็วคงตัวไม่ได้แปลว่าความเร่งเป็นศูนย์
    \item อธิบายทิศของความเร็วและความเร่งในการเคลื่อนที่เป็นวงกลมได้
    \item ใช้สูตรความเร่งสู่ศูนย์กลาง $a_c=\frac{v^2}{r}$ และ $a_c=\omega^2r$ ได้
    \item เข้าใจความสัมพันธ์ระหว่างคาบ ความถี่ อัตราเร็วเชิงมุม และอัตราเร็วเชิงเส้น
    \item ใช้กฎข้อที่สองของนิวตันในแนวรัศมีได้
    \item เข้าใจว่าแรงสู่ศูนย์กลางไม่ใช่แรงชนิดใหม่ แต่เป็นแรงลัพธ์ในแนวรัศมี
    \item วิเคราะห์โจทย์ลูกตุ้มกรวย รถเลี้ยว ทางโค้งเอียง และวงกลมแนวดิ่งได้
\end{enumerate}

\begin{tcolorbox}[title=แนวคิดสำคัญของบทนี้]
การเคลื่อนที่เป็นวงกลมต้องมีความเร่งเข้าหาศูนย์กลางเสมอ แม้อัตราเร็วจะคงตัวก็ตาม จุดสำคัญของบทนี้คือการใช้กฎของนิวตันในแนวรัศมี
\[
\sum F_r=m\frac{v^2}{r}
\]
โดยต้องจำไว้ว่าแรงสู่ศูนย์กลางไม่ใช่แรงชนิดใหม่ แต่คือแรงลัพธ์ที่ชี้เข้าสู่ศูนย์กลาง
\end{tcolorbox}

\newpage{}

% =========================================================
\section{การเคลื่อนที่เป็นวงกลมสม่ำเสมอ}

\subsection{อัตราเร็วคงตัวแต่ความเร่งไม่เป็นศูนย์}

ในการเคลื่อนที่เป็นวงกลมสม่ำเสมอ วัตถุเคลื่อนที่ด้วยอัตราเร็วคงตัว

\[
v=\text{constant}
\]

แต่เวกเตอร์ความเร็วไม่ได้คงตัว เพราะทิศของความเร็วเปลี่ยนตลอดเวลา

ความเร็วของวัตถุที่เคลื่อนที่บนวงกลมจะชี้ตามแนวสัมผัสของวงกลมเสมอ

\[
\vec{v} \text{ ชี้ตามแนวสัมผัส}
\]

ส่วนความเร่งของการเคลื่อนที่เป็นวงกลมสม่ำเสมอจะชี้เข้าหาศูนย์กลางวงกลม

\[
\vec{a}_c \text{ ชี้เข้าหาศูนย์กลาง}
\]

\begin{tcolorbox}[title=ข้อควรจำ]
อัตราเร็วเป็นขนาดของความเร็ว ถ้าอัตราเร็วคงตัว แต่ทิศของความเร็วเปลี่ยน เวกเตอร์ความเร็วยังเปลี่ยนอยู่ ดังนั้นวัตถุยังมีความเร่ง
\end{tcolorbox}

\subsection{ความเร่งสู่ศูนย์กลาง}

ความเร่งที่ชี้เข้าหาศูนย์กลางวงกลมเรียกว่า ความเร่งสู่ศูนย์กลาง

\[
a_c=\frac{v^2}{r}
\]

โดยที่

\[
v=\text{อัตราเร็วเชิงเส้น}
\]

และ

\[
r=\text{รัศมีของวงกลม}
\]

\begin{tcolorbox}[title=ความเร่งสู่ศูนย์กลาง]
\[
a_c=\frac{v^2}{r}
\]

ทิศของ $\vec{a}_c$ ชี้เข้าหาศูนย์กลางวงกลม
\end{tcolorbox}

% =========================================================
\section{คาบ ความถี่ และอัตราเร็วเชิงมุม}

\subsection{คาบและความถี่}

คาบ $T$ คือเวลาที่วัตถุใช้เคลื่อนที่ครบหนึ่งรอบ

\[
T=\text{เวลาต่อหนึ่งรอบ}
\]

ความถี่ $f$ คือจำนวนรอบต่อหนึ่งวินาที

\[
f=\frac{1}{T}
\]

ถ้าวัตถุเคลื่อนที่เป็นวงกลมรัศมี $r$ ระยะทางหนึ่งรอบคือ

\[
2\pi r
\]

ดังนั้นอัตราเร็วเชิงเส้นคือ

\[
v=\frac{2\pi r}{T}
\]

\subsection{อัตราเร็วเชิงมุม}

อัตราเร็วเชิงมุม $\omega$ คือมุมที่กวาดได้ต่อเวลา

\[
\omega=\frac{\Delta\theta}{\Delta t}
\]

ถ้าเคลื่อนที่ครบหนึ่งรอบ มุมที่กวาดได้คือ

\[
2\pi
\]

เรเดียน ดังนั้น

\[
\omega=\frac{2\pi}{T}
\]

และเพราะ

\[
v=\frac{2\pi r}{T}
\]

จึงได้

\[
v=\omega r
\]

แทนลงใน

\[
a_c=\frac{v^2}{r}
\]

จะได้

\[
a_c=\frac{(\omega r)^2}{r}
\]

ดังนั้น

\[
a_c=\omega^2r
\]

\begin{tcolorbox}[title=สูตรพื้นฐานของการเคลื่อนที่เป็นวงกลม]
\[
f=\frac{1}{T}
\]

\[
v=\frac{2\pi r}{T}
\]

\[
\omega=\frac{2\pi}{T}
\]

\[
v=\omega r
\]

\[
a_c=\frac{v^2}{r}=\omega^2r
\]
\end{tcolorbox}

% =========================================================
\section{แรงสู่ศูนย์กลาง}

\subsection{แรงสู่ศูนย์กลางไม่ใช่แรงชนิดใหม่}

ในการเคลื่อนที่เป็นวงกลม วัตถุต้องมีความเร่งเข้าหาศูนย์กลาง ดังนั้นจากกฎข้อที่สองของนิวตัน ต้องมีแรงลัพธ์เข้าหาศูนย์กลางด้วย

\[
\sum F_r=ma_c
\]

แทน

\[
a_c=\frac{v^2}{r}
\]

จะได้

\[
\sum F_r=m\frac{v^2}{r}
\]

\begin{tcolorbox}[title=แนวคิดสำคัญ]
แรงสู่ศูนย์กลางไม่ใช่แรงชนิดใหม่ แต่คือแรงลัพธ์ในแนวรัศมี

แรงจริงที่ทำหน้าที่เป็นแรงสู่ศูนย์กลางอาจเป็นแรงตึงเชือก แรงปกติ แรงเสียดทาน หรือแรงโน้มถ่วง ขึ้นกับสถานการณ์
\end{tcolorbox}

\subsection{วิธีตั้งสมการ}

เลือกแกนรัศมีให้ชี้เข้าหาศูนย์กลางวงกลม แล้วเขียน

\[
\sum F_r=m\frac{v^2}{r}
\]

ถ้าอัตราเร็วเปลี่ยนด้วย ให้มีแกนสัมผัสเพิ่มอีกแกนหนึ่ง

\[
\sum F_t=ma_t
\]

โดย

\[
a_t=\frac{dv}{dt}
\]

\begin{tcolorbox}[title=ข้อควรระวัง]
อย่าวาดแรง $m\frac{v^2}{r}$ ลงใน free-body diagram เพราะไม่ใช่แรงจริง แต่เป็นผลของแรงลัพธ์ในแนวรัศมี
\end{tcolorbox}

% =========================================================
\section{ลูกตุ้มกรวย}

ลูกตุ้มกรวยคือมวล $m$ ผูกกับเชือกแล้วเคลื่อนที่เป็นวงกลมในแนวระดับ โดยเชือกทำมุม $\theta$ กับแนวดิ่ง

แรงที่กระทำต่อมวลมีสองแรง

\[
mg \text{ ลง}
\]

และ

\[
T \text{ ตามแนวเชือก}
\]

แตกแรงตึงเชือกได้เป็น

\[
T\cos\theta \text{ ขึ้น}
\]

และ

\[
T\sin\theta \text{ เข้าหาศูนย์กลาง}
\]

แนวดิ่งไม่มีความเร่ง

\[
T\cos\theta=mg
\]

แนวรัศมีมีความเร่งสู่ศูนย์กลาง

\[
T\sin\theta=m\frac{v^2}{r}
\]

หารสมการแนวรัศมีด้วยสมการแนวดิ่ง

\[
\frac{T\sin\theta}{T\cos\theta}
=
\frac{m\frac{v^2}{r}}{mg}
\]

จึงได้

\[
\tan\theta=\frac{v^2}{rg}
\]

ดังนั้น

\[
v=\sqrt{rg\tan\theta}
\]

\begin{tcolorbox}[title=ลูกตุ้มกรวย]
\[
T\cos\theta=mg
\]

\[
T\sin\theta=m\frac{v^2}{r}
\]

\[
v=\sqrt{rg\tan\theta}
\]
\end{tcolorbox}

% =========================================================
\section{รถเลี้ยวบนถนนราบ}

รถเลี้ยวบนถนนราบต้องมีแรงเข้าหาศูนย์กลาง ถ้ารถไม่ไถล แรงที่ทำหน้าที่เป็นแรงสู่ศูนย์กลางคือแรงเสียดทานสถิตระหว่างยางกับถนน

แรงเสียดทานสถิตสูงสุดคือ

\[
f_{s,\max}=\mu_sN
\]

บนถนนราบ

\[
N=mg
\]

ดังนั้น

\[
f_{s,\max}=\mu_smg
\]

เพื่อให้รถเลี้ยวได้โดยไม่ไถล ต้องมี

\[
m\frac{v^2}{r}\leq \mu_smg
\]

ตัด $m$

\[
\frac{v^2}{r}\leq \mu_sg
\]

ดังนั้น

\[
v_{\max}=\sqrt{\mu_sgr}
\]

\begin{tcolorbox}[title=รถเลี้ยวบนถนนราบ]
\[
v_{\max}=\sqrt{\mu_sgr}
\]
\end{tcolorbox}

% =========================================================
\section{ทางโค้งเอียงไร้แรงเสียดทาน}

ถนนโค้งเอียงทำให้แรงปกติมีองค์ประกอบเข้าหาศูนย์กลางได้ แม้ไม่มีแรงเสียดทาน

ให้ถนนเอียงทำมุม $\theta$ กับแนวระดับ และรถเคลื่อนที่เป็นวงกลมรัศมี $r$

แรงที่กระทำต่อรถมี

\[
mg \text{ ลง}
\]

และ

\[
N \text{ ตั้งฉากกับถนน}
\]

แตกแรงปกติ

\[
N\cos\theta \text{ ขึ้น}
\]

\[
N\sin\theta \text{ เข้าหาศูนย์กลาง}
\]

แนวดิ่งไม่มีความเร่ง

\[
N\cos\theta=mg
\]

แนวรัศมี

\[
N\sin\theta=m\frac{v^2}{r}
\]

หารสองสมการ

\[
\tan\theta=\frac{v^2}{rg}
\]

ดังนั้น

\[
v=\sqrt{rg\tan\theta}
\]

\begin{tcolorbox}[title=ทางโค้งเอียงไร้แรงเสียดทาน]
\[
\tan\theta=\frac{v^2}{rg}
\]

\[
v=\sqrt{rg\tan\theta}
\]
\end{tcolorbox}

% =========================================================
\section{วงกลมแนวดิ่ง}

ในการเคลื่อนที่เป็นวงกลมแนวดิ่ง ทิศเข้าสู่ศูนย์กลางเปลี่ยนไปตามตำแหน่ง และน้ำหนัก $mg$ ชี้ลงเสมอ

\subsection{ตำแหน่งล่างสุด}

ที่ตำแหน่งล่างสุด ทิศเข้าสู่ศูนย์กลางคือขึ้น

แรงในแนวรัศมีมี

\[
T \text{ ขึ้น}
\]

และ

\[
mg \text{ ลง}
\]

ใช้

\[
\sum F_r=m\frac{v^2}{r}
\]

จึงได้

\[
T-mg=m\frac{v^2}{r}
\]

ดังนั้น

\[
T=mg+m\frac{v^2}{r}
\]

\subsection{ตำแหน่งบนสุด}

ที่ตำแหน่งบนสุด ทิศเข้าสู่ศูนย์กลางคือลง

แรงในแนวรัศมีมี

\[
T \text{ ลง}
\]

และ

\[
mg \text{ ลง}
\]

ดังนั้น

\[
T+mg=m\frac{v^2}{r}
\]

\[
T=m\frac{v^2}{r}-mg
\]

ถ้าเชือกเพิ่งจะไม่หย่อนที่ตำแหน่งบนสุด

\[
T=0
\]

ดังนั้น

\[
mg=m\frac{v^2}{r}
\]

จึงได้

\[
v_{\min}=\sqrt{gr}
\]

\begin{tcolorbox}[title=วงกลมแนวดิ่ง]
ตำแหน่งล่างสุด

\[
T=mg+m\frac{v^2}{r}
\]

ตำแหน่งบนสุด

\[
T=m\frac{v^2}{r}-mg
\]

เงื่อนไขเชือกไม่หย่อนที่บนสุด

\[
v_{\min}=\sqrt{gr}
\]
\end{tcolorbox}

% =========================================================
\section{สรุปสูตรสำคัญ}

\begin{tcolorbox}[title=สรุปบท: การเคลื่อนที่เป็นวงกลม]
\[
f=\frac{1}{T}
\]

\[
v=\frac{2\pi r}{T}
\]

\[
\omega=\frac{2\pi}{T}
\]

\[
v=\omega r
\]

\[
a_c=\frac{v^2}{r}=\omega^2r
\]

\[
\sum F_r=m\frac{v^2}{r}
\]

\[
a_t=\frac{dv}{dt}
\]

\[
a=\sqrt{a_r^2+a_t^2}
\]
\end{tcolorbox}

\begin{tcolorbox}[title=สรุปสูตรประยุกต์]
ลูกตุ้มกรวย

\[
v=\sqrt{rg\tan\theta}
\]

รถเลี้ยวบนถนนราบ

\[
v_{\max}=\sqrt{\mu_sgr}
\]

ทางโค้งเอียงไร้แรงเสียดทาน

\[
v=\sqrt{rg\tan\theta}
\]

วงกลมแนวดิ่งตำแหน่งล่างสุด

\[
T=mg+m\frac{v^2}{r}
\]

วงกลมแนวดิ่งตำแหน่งบนสุด

\[
T=m\frac{v^2}{r}-mg
\]

เงื่อนไขเชือกไม่หย่อนที่บนสุด

\[
v_{\min}=\sqrt{gr}
\]
\end{tcolorbox}

% =========================================================
\section{ข้อสอบจริง สอวน. ที่ใช้ความรู้บทการเคลื่อนที่เป็นวงกลม}

% =========================================================
\subsection{ปี 2562 ข้อ 6: มวลบนผิวกรวยลื่นเคลื่อนที่เป็นวงกลม}

\vspace{1.5em}

\IfFileExists{img/posn62q6.png}{
\begin{center}
    \includegraphics[width=1\linewidth]{img/posn62q6.png}
\end{center}
}{
\begin{tcolorbox}[title=ตำแหน่งภาพที่แนะนำ]
ให้ตัดภาพข้อสอบปี 2562 ข้อ 6 แล้วบันทึกเป็น

\[
\texttt{img/posn62q6.png}
\]
\end{tcolorbox}
}

\vspace{1em}

\begin{examplebox}{เฉลยละเอียด}
มวลเล็กเคลื่อนที่เป็นวงกลมในแนวระดับบนผิวด้านในของกรวยลื่น ที่ความสูง

\[
h
\]

จากพื้น

จากรูป กรวยทำให้รัศมีของวงกลมที่ระดับความสูงนี้เป็น

\[
r=h
\]

แรงที่กระทำต่อมวลมีสองแรง

\[
mg \text{ ลง}
\]

และแรงปกติจากผิวกรวย

\[
N
\]

เพราะผิวกรวยลื่น จึงไม่มีแรงเสียดทาน

มวลเคลื่อนที่เป็นวงกลมในแนวระดับ ดังนั้นแนวดิ่งไม่มีความเร่ง

\[
a_y=0
\]

แต่แนวรัศมีมีความเร่งสู่ศูนย์กลาง

\[
a_c=\frac{v^2}{r}
\]

จากรูป แนวแรงปกติทำมุมที่ทำให้องค์ประกอบแนวดิ่งและแนวรัศมีเท่ากัน ดังนั้น

\[
N_y=N_r
\]

พิจารณาแนวดิ่ง

\[
N_y=mg
\]

ดังนั้น

\[
N_r=mg
\]

ใช้กฎข้อที่สองของนิวตันในแนวรัศมี

\[
N_r=m\frac{v^2}{r}
\]

แทนค่า

\[
mg=m\frac{v^2}{r}
\]

ตัด $m$

\[
g=\frac{v^2}{r}
\]

ดังนั้น

\[
v^2=gr
\]

และเพราะ

\[
r=h
\]

จึงได้

\[
v^2=gh
\]

ดังนั้น

\[
\boxed{v=\sqrt{gh}}
\]
\end{examplebox}